\section{Introduction}

Solving large sparse systems of linear equations is the central task of
scientific computing, and electromagnetics is no exception: discretizing
Maxwell's equations with finite differences or finite elements turns the
analysis and design of a photonic device into a sparse linear system with
millions to billions of unknowns~\citep{yee1966,Oskooi2010,ShinFan2012}.
Optical analog computing runs this correspondence in the opposite direction.
Instead of computing the response of a given structure, one designs a structure
whose response computes something else -- e.g., a derivative, a convolution or
the solution of an equation~\citep{Silva2014,ZangenehNejad2021,McMahon2023}.
Metamaterials and inverse-designed metastructures have been shown to solve
integral equations~\citep{estakhri2019,Camacho2021,cordaro2023}, reconfigurable
networks invert matrices~\citep{tzarouchis2022,chen2024}, and integrated
photonic circuits solve differential equations~\citep{yang2014,ye2024}. These
build on a broader effort to perform linear algebra with light, from universal
interferometer meshes~\citep{reck1994,Miller2013SelfConfiguring,Carolan2015,Clements2016,Bogaerts2020}
to optical neural networks and photonic tensor
cores~\citep{Shen2017,Lin2018,Feldmann2021,Wetzstein2020,Shastri2021,Wright2022}.

Mathematical constructions also establish the expressive power of
electromagnetic networks. \citet{MiltonSeppecher2010} characterize
terminal-response matrices at a fixed frequency, including arbitrary real
symmetric matrices in the lossless case. Their continuum construction uses
dielectric and magnetic elements with parameters tending to zero or infinity,
and leaves the accuracy of approximations with finite material parameters
open. Reconfigurable wave networks support general matrix inversion,
including generalized inverses, with bounds on convergence
time~\citep{tzarouchis2022}.

Whether such devices can ever offer a computational advantage is decided by
details that are easy to overlook in a demonstration: the cost of programming
the device, the precision with which it must be built and read out, the
sensitivity of the answer to imperfections, and the size of the structure as
the problem grows~\citep{Aaronson2015FinePrint,McMahon2023}. The same questions
were asked of analog computers long before the current interest in photonics,
and they have a precise formulation in the theory of
computation~\citep{Shannon1941,Vergis1986,BSS1989,SiegelmannSontag1994,BournezPouly2021}.
Physical systems as simple as waves, rays and fluids can embed universal
computation~\citep{PourElRichards1981,Moore1990,ReifTygarYoshida1994,Cubitt2015,Tao2016,Cardona2021},
but these constructions establish what is computable in principle, not what a
physical system computes efficiently. For optical equation solvers, the
relevant question is quantitative: can general sparse systems be represented
by a passive dielectric with fixed material bounds while controlling
construction size, conditioning and precision?

Electromagnetic simulation faces the same question from the other side. The
linear systems produced by Maxwell solvers are highly structured -- local,
three-dimensional and determined by a scalar material distribution -- and fast
solvers exploit this structure wherever they
can~\citep{George1973,LiptonRoseTarjan1979,EngquistYing2011CPAM,Martinsson2019}.
However, time-harmonic wave problems are notoriously hard to solve
iteratively~\citep{ErnstGander2012,GanderZhang2019,MoiolaSpence2014}, and it
is not known whether their structure makes them fundamentally easier than
general sparse systems. For other structured classes this question has sharp
answers. Graph Laplacians can be solved in nearly linear
time~\citep{SpielmanTeng2014}, whereas the stiffness matrices of trusses and
multi-commodity flow problems are as hard as general sparse
systems~\citep{kyngzhang2017,DingKyngProbstGutenbergZhang2022,DingKyngZhang2022}.
To the best of our knowledge, no reduction from general sparse systems to
continuum Maxwell equations with fixed permeability and bounded scalar
permittivity has established explicit bounds on size, conditioning and
precision.

Both questions ask how expressive a dielectric is, i.e., which equations can be
written into a material distribution, and at what cost. Here, we answer this
question with an explicit compiler. Given a sparse system $Mx=b$ with $m$
nonzero coefficients, polynomial condition number and coefficients of
$O(\log m)$ bits, the compiler outputs a three-dimensional dielectric inside a
perfectly conducting box, a source and a finite list of field measurements,
such that every field approximately satisfying Maxwell's equations yields an
approximate solution of $Mx=b$. The matrix may be nonsymmetric, its entries may
have either sign and its sparsity pattern is arbitrary, while the dielectric is
lossless, isotropic and has a permittivity between 1 and 25. The volume of the
structure and the computational cost of encoding and decoding all scale as
$\Ot(m^{3/2})$, where the tilde hides logarithmic factors. We prove this for
the continuum Maxwell equations and give a computer-assisted construction with
the same scaling for the Yee discretization~\citep{yee1966} used by
finite-difference solvers.

Physically, the construction is the familiar picture of coupled defect modes
in a photonic crystal~\citep{Yariv1999,Busch2003,Joannopoulos2008}, made
quantitative. Each unknown is carried by the amplitude of a localized mode,
couplings between unknowns arise from the overlap of neighboring modes, and
auxiliary modes mediate the signed and nonsymmetric couplings that a
reciprocal dielectric can not provide directly. The proof controls everything
this picture neglects -- the remaining field components, the unwanted
couplings summed over the whole device, the finite boundary and near-resonances
of the cavity -- tightly enough that the encoded equations survive. Read in
reverse, the compiler turns any electromagnetic solver into a sparse linear
solver. A solver that handles the resulting family of structures in time
$N^\theta$, with $N$ the size of the structure, therefore solves sparse linear
systems in time $\Ot(m^{3/2}+m^{3\theta/2})$.

We first describe the compiler and its guarantees, then the physical mechanism
that makes it work and the recovery of answers from approximate fields. We then
discuss what the result implies for optical computing and for electromagnetic
simulation. Methods summarizes the formal statements, and the \SI\ contains the
complete proofs.
